\nuevaReceta{Berenjenas con Miel de Caña}{30 min}{receta:berenjenas_miel}

    \begin{minipage}[t]{0.35\textwidth}
        \section*{Ingredientes}
        \begin{itemize}[leftmargin=*]
            \item 2 Berenjenas
            \item Harina
            \item Sal
            \item Aceite para freír
            \item Miel de caña
        \end{itemize}
    \end{minipage}
    \hfill
    \begin{minipage}[t]{0.6\textwidth}
        \section*{Preparación}
        \begin{enumerate}
            \item \textbf{Preparar las berenjenas:} Lavar las berenjenas y cortarlas en bastones finos.
            \item \textbf{Quitar el amargor:} Ponerlas en agua con un poco de sal durante un par de horas para que suelten agua. Después, secarlas bien con papel de cocina.
            \item \textbf{Enharinar:} Pasar las berenjenas por una mezcla de cerveza y harina, sacudiendo el exceso para que no queden apelmazadas.
            \item \textbf{Freír:} Freírlas en abundante aceite caliente hasta que queden doradas y crujientes.
            \item \textbf{Escurrir:} Sacarlas a un plato con papel absorbente para retirar el exceso de aceite.
            \item \textbf{Servir:} Colocarlas en una fuente y echar por encima un hilo de miel de caña justo antes de servir.
        \end{enumerate}
    \end{minipage}

    \vspace{1em}
    \begin{tcolorbox}[colback=orange!5, colframe=orange!50, title=Consejo, size=small]
        Si quieres un contraste más fino, sirve la miel de caña templada y échala al final para que las berenjenas mantengan mejor el crujiente.
    \end{tcolorbox}
\end{receta}
